\section{Conclusiones}\label{sec:conclusiones}

La entrevista identifica como prioridad organizar el registro y seguimiento de los equipos.
El caso de una reparación detenida por un repuesto muestra la dificultad de reconstruir su
estado entre varias anotaciones. Esto respalda una propuesta centrada en localizar servicios
y consultar su avance, manteniendo la atención directa del taller.

El alcance distingue registro de reparaciones y cobros de funciones más amplias, como inventario,
mensajería automática o facturación. Estas inclusiones y exclusiones son propuestas del equipo
y deben confirmarse con el responsable.

Se cuenta con consentimiento, transcripción y audio registrados en el manifiesto. La entrevista
aporta testimonio, pero no una medición de frecuencia, tiempos o pagos omitidos. El siguiente
paso es revisar ejemplos anonimizados, acordar metas y validar la propuesta; el cuestionario y
la observación están preparados, todavía sin aplicar.
